% vim: textwidth=78 ai spell spelllang=en filetype=tex sw=2

\chapter{Participating SV-COMP}\label{sec:svcomp}
We had a couple of tools, we were checking the Linux kernel and finding real
errors. We have not considered any other use our tools until we were pointed
by the authors of \textsc{Predator}~\cite{predator} to the past
\emph{Competition on Software Verification} (SV-COMP) 2012~\cite{sv-comp12} in
which they participated. We decided to participate in the next year's
competition with our tool \symbio~\cite{symbiotic}. And since one of authors
of every tool is asked to be in the competition program committee, the author
of this thesis willingly took the role.

We have already provided a detailed description of the \symbio internals in
Chapter~\ref{sec:fmics}. Now we only sum up that the technique implemented in
\symbio combines instrumentation, program slicing, and symbolic execution in
order to detect bugs described by automata. More precisely, we instrument a
given program with code that tracks runs of automata representing various
erroneous behaviors. If an instrumented automaton enters an error location
during a program execution, then the original program contains a bug specified
by the automaton. After instrumentation, we slice~\cite{slicing} the program
to reduce its size without affecting runs of the instrumented automata.
Finally, we symbolically execute~\cite{se_king} the sliced program to find
errors in the program.

As reachability of an \code{ERROR} label is the only bug considered in the
SV-COMP, we have modified our instrumentor to put \code{assert(0)} function
calls at \code{ERROR} labels in the code. Given the instrumented code, we
execute \clang to produce \llvm bytecode. This is in turn inter-procedurally
sliced with respect to slicing criteria, which are the instrumented
\code{assert} calls. In other words we remove all the code except that has an
effect on reachability of \code{assert} calls. The sliced \llvm bitcode is
finally symbolically executed by \klee~\cite{klee}. There are several possible
outputs that \klee may generate. It can either find a reachable \code{assert}
and report it, or finish the computation without any report, or terminate in
some errant way (out of time, out of memory, invalid memory dereference, or
some internal error for example). We map these to the demanded answers:
\code{UNSAFE}, \code{SAFE}, or \code{UNKNOWN} respectively.  This is taken
care of in a simple scripted filter.

Participation in this competition is documented in a paper published in
Springer's LNCS~\cite{symbiotic-svcomp} and presented on the competition as a
part of the 19th International Conference on Tools and Algorithms for the
Construction and Analysis of Systems. There is also a competition report from
2013 by Beyer~\cite{sv-comp13}.

%}}}
%{{{ Software architecture

\section{Software Architecture}
\paragraph{Libraries/External Tools} For the code translation and for
operations with the \textsc{Llvm} bitcode we use
\textsc{Llvm}/\textsc{Clang} 3.1\footnote{\url{http://llvm.org}}. The
symbolic executor is \textsc{Klee} which itself uses the
\textsc{Stp} constraint solver~\cite{stp}. We obtained both \textsc{Klee} and
\textsc{Stp} from the respective GIT repositories and used the snapshots. For
more information about \textsc{Klee}, see~\cite{klee}.

\paragraph{Software Structure and Architecture} The architecture described
in the introduction above is summarized in Figure~\ref{fig:pipeline}.  The
slicer is written as a plug-in for the optimizer \texttt{opt} from the
\textsc{Llvm} suite. It is publicly available in a separate repository at
\url{https://github.com/jirislaby/LLVMSlicer/}. The slicer improves the
symbolic execution considerably. For ease of use, all parts of the tool
pipeline are one by one run by a single script \texttt{runme}.

\begin{figure}[tb]
\begin{center}
  \tikzstyle{stat} = [thick,draw,minimum size=1cm,node distance=1.5cm]
  \tikzstyle{pre} = [<-,shorten <=1pt,>=stealth',semithick]
  \begin{tikzpicture}[font=\sffamily]
    \node (START) {};
    \node [stat,right=of START] (INSTR) {Instrumentor}
      edge [pre] node [label=above:{C}] {} (START);
    \node [stat,right=of INSTR] (CLANG) {\textsc{Clang}}
      edge [pre] node [label=above:{C}] {} (INSTR);
    \node [stat,right=2cm of CLANG] (SLICER) {Slicer}
      edge [pre] node [label=above:{LLVM}] {} (CLANG);
    \node [stat,below=9mm of SLICER] (KLEE) {\textsc{Klee}}
      edge [pre] node [label=right:{LLVM}] {} (SLICER);
    \node [stat,left=2.8cm of KLEE] (FILTER) {Filter}
      edge [pre] node [label=above:{\textsc{Klee} result}] {} (KLEE);
    \node [left=4.9cm of FILTER] (FINISH) {}
      edge [pre] node [label=above:{\texttt{SAFE/UNSAFE/UNKNOWN}}] {} (FILTER);
  \end{tikzpicture}
\end{center}
  \caption{Pipeline of the tool}
\label{fig:pipeline}
\end{figure}

\paragraph{Implementation Technology} The instrumentation is performed by a
\texttt{bash} script using \texttt{sed}. The final filter also uses
\texttt{bash} with the help of \texttt{grep}. The rest of the toolchain is
written in C++ and compiled using \texttt{gcc}.

%}}}
%{{{ Discussion

\section{Discussion of Strengths and Weaknesses of the Approach}

The strength of the tool lies in its high precision of the answers. In theory,
the only source of incorrect answers is the slicer: it can completely remove
an infinite loop in some cases and thus an unreachable \texttt{ERROR} label
located below the loop may become reachable. However, there is no such case
in the competition benchmarks. 

All incorrect answers produced by our tool in the competition are due to bugs
in implementation.  Since the tool submission, we have fixed most of the bugs
and improved the implementation a bit. The biggest change on the competition
benchmarks can be seen in the category \emph{FeatureChecks} (118 files):

\begin{center}
  \setlength{\tabcolsep}{3pt}
  \begin{tabular}{|l|c|c|}
    \hline
    & Competition Version & Current Version \\ \hline 
    Correct Answers \texttt{SAFE/UNSAFE} & 81 & 116 \\ \hline
    Incorrect Answers \texttt{SAFE/UNSAFE} & 17 & 0 \\ \hline
    \texttt{UNKNOWN} (including timeouts) & 20 & 2 \\ \hline
  \end{tabular}
\end{center}

In general, the main weakness of the tool is a high percentage of
\texttt{UNKNOWN} results. These results come mainly from high computation
cost of the symbolic execution of programs with loops or recursion. This
problem is relieved by slicing, but there are still many cases where the
sliced code remains complex and symbolic execution runs out of time or
memory. Other sources of \texttt{UNKNOWN} results are internal errors of
\textsc{Klee} and general limitations of constraint solving.

%}}}
%{{{ Tool setup and configuration

\section{Tool Setup and Configuration}
\paragraph{Download and Installation Instructions}
\begin{itemize}
\item \emph{Requirements:} \texttt{llvm-3.1}, \texttt{clang-3.1}.
\item Download \symbio 1 at: \url{https://sf.net/projects/symbiotic/}.
\item Change the current directory to \texttt{/opt} (this location is
	needed for \textsc{Klee}).
\item Untar the archive with the tool.
\item Change directory into \texttt{/opt/symbiotic}.
\item Run \texttt{./runme <benchmark.c>} for each source file in the set.
\end{itemize}

\paragraph{Results Reported} \texttt{SAFE/UNSAFE/UNKWNOWN} answers match
the competition rules. The counterexample for an error in some
\texttt{<benchmark.c>} is generated for each error path in the benchmark to
\texttt{<benchmark.c>-klee-out/}. There is also a sliced code referred by all
the error paths.

%}}}
%{{{ Software project and contributors

\section{Software Project and Contributors}
The concept of the tool has been developed by the authors of this paper. The
tool was implemented by Jiri Slaby (contact person), Marek Trt\'ik, and Ben
Liblit (several fixes).  It is available under the GNU GPLv2 License and is
hosted by the Faculty of Informatics, Masaryk University.

%}}}
